\scanpage{58}
Automorphismes du $n$-cube.

Considérons le $n$-cube de maquette $(B,\sigma)$, soit $I=B/\sigma$. Un automorphisme correspond à un autom. \inserted{$u$} de la maquette, induisant un autom $u_I$ de $I$, d'où une décomposition de $I$ en orbites
\[
 I=\coprod_{\deleted{\unsure{i}}\alpha\deleted{\unsure{\in A}}}I_\alpha
\]
($u_I$ est connu modulo conjugaison quand on connaît les cardinaux des $I_\alpha$, rangés par ordre croissant \ldots)

Soit $B_\alpha$ l'image inverse de $I_\alpha$, donc $B_\alpha$ est stable par $u$, \deleted{donc} \deleted{réunion d'orbites, \unsure{il y en a|dont|il y a des} chacune s'envoie \uline{sur} $I_\alpha$ avec un degré déterminé. Donc 2 cas à envisager $1^\circ$) $B_\alpha$ est une orbite. Soit $n_\alpha$} \inserted{la situation est $\sim$ somme des $(B_\alpha,u_\alpha=u|B_\alpha)$ sur $I_\alpha=B_\alpha/(\underbrace{\sigma|B_\alpha}_{\unsure{\sigma_\alpha}})$,} et on est amené à étudier chacun de ses \unsure{c.|cas} (Géométriquement, $C=\prod C_\alpha$, où $C_\alpha$ est le cube associé à $(B_\alpha,\unsure{\sigma_\alpha|\sigma}\deleted{\unsure{I_\alpha|\mid B_\alpha}}$, et $u=\prod u_\alpha$,\ldots). Celà nous ramène à étudier le cas où $u_I$ est transitif sur $I$.

$B$ est réunion d'orbites, dont chacune s'envoie sur $I$ (l'unique orbite \deleted{\unsure{$I_\alpha$}} de $u_I$ dans $I$), avec un degré déterminé. Donc il y a 2 cas

a) Deux orbites isomorphes à $I$ \deleted{\unsure{à lui|$\xrightarrow{\sim}$}} par $p:B\unsure{-|\to} I$. Alors, si \deleted{\unsure{$i$}} $\unsure{E|\varepsilon}$ est l'ens des 2 orbites, on a canoniquement
\[
 B\simeq E\times I,\qquad u=\mathrm{Id}_E\times\deleted{\unsure{\mathrm{Id}_I|\mathrm{Id}}}u_I
\]
% archive sheet number: 57
